\documentclass[11pt,a4paper]{article}

% Packages
\usepackage[utf8]{inputenc}
\usepackage[T1]{fontenc}
\usepackage[margin=0.75in]{geometry}
\usepackage{enumitem}
\usepackage{hyperref}
\usepackage{titlesec}
\usepackage{xcolor}

% Formatting
\pagestyle{empty}
\setlength{\parindent}{0pt}
\setlist[itemize]{leftmargin=*, noitemsep, topsep=3pt}

% Section formatting
\titleformat{\section}{\Large\bfseries}{}{0em}{}[\titlerule]
\titlespacing*{\section}{0pt}{12pt}{6pt}

\titleformat{\subsection}[runin]{\bfseries}{}{0em}{}
\titlespacing*{\subsection}{0pt}{8pt}{0pt}

% Colors
\definecolor{linkcolor}{RGB}{0,102,204}
\hypersetup{
    colorlinks=true,
    urlcolor=linkcolor
}

\begin{document}

% Header
\begin{center}
    {\LARGE \textbf{ERIC MAYEUX}}\\
    \vspace{4pt}
    {\large Chargé de projet | Project Manager}\\
    \vspace{4pt}
    Montréal, QC, Canada\\
    \href{mailto:mayeux.e@gmail.com}{mayeux.e@gmail.com} | 438-884-7645
\end{center}

\vspace{8pt}

\section*{Profil Professionnel}
\noindent
\textbf{Chargé de projet} avec 10+ ans d'expérience en \textbf{planification et coordination de projets} TI et opérationnels. Expert en \textbf{mobilisation des parties prenantes}, suivi des livrables et \textbf{gestion des risques et enjeux}. Maîtrise des outils collaboratifs (Jira, MS Project) et des méthodologies \textbf{Agile (Scrum, Kanban)} et Waterfall. Certifié \textbf{PMP} (en cours) et CSM. Bilinguisme français/anglais. Forte orientation vers l'\textbf{amélioration continue} et l'efficacité opérationnelle.

\section*{Compétences Techniques}

\textbf{Gestion de projet:} Planification et cadrage de projets, suivi des livrables, échéanciers, budgets, amélioration continue, bureau de projets (PMO)

\textbf{Mobilisation \& Communication:} Parties prenantes internes/externes, animation de réunions et ateliers, communication fluide, leadership d'influence

\textbf{Gestion des risques:} Identification, registre des risques et enjeux, plans de mitigation, support à la prise de décision

\textbf{Documentation:} Charte de projet, plan de projet, registre des risques, rapport d'avancement, bilan des leçons apprises

\textbf{Outils:} Jira, MS Project, Confluence, Datadog, Splunk, tableaux de bord KPI, Asana, Trello

\textbf{Méthodologies:} Agile, Scrum, SAFe, Kanban, Waterfall, PMP

\textbf{Intelligence Artificielle:} Agents IA, automatisation de tâches de gestion et d'analyse

\vspace{6pt}

\section*{Réalisations}
\begin{itemize}
    \item \textbf{Planification et coordination} de releases multi-équipes à la BNC : de l'estimation à la mise en production
    \item \textbf{Mobilisation des parties prenantes} dans un contexte SAFe : gouvernance, communication, alignement
    \item Mise en place de mécanismes de \textbf{suivi de l'avancement} : tableaux de bord KPI, rapports d'avancement
    \item \textbf{Amélioration continue} : audit qualité, état des lieux, optimisation des pratiques de gestion
    \item Gestion de \textbf{plusieurs projets en parallèle} dans un contexte bancaire agile (SAFe Release Train)
    \item Mise en place d'agents IA pour l'automatisation de tâches de \textbf{gestion et de documentation}
\end{itemize}

\section*{Expérience Professionnelle}

\subsection*{Scrum Master -- Gestion de projet | Project Manager} \hfill \textit{Novembre 2025 -- Présent}\\
\textbf{Banque Nationale du Canada (BNC)} -- Montréal, QC\\
Coordination de projets TI stratégiques : planification, suivi des livrables, mobilisation des équipes.
\begin{itemize}
    \item Encadre des équipes IT multidisciplinaires (QA, Dev, PO, Analystes) vers les objectifs de livraison
    \item \textbf{Pilote les releases} : cadrage, planification, coordination inter-équipes, suivi des dépendances
    \item Identifie et gère les \textbf{risques et enjeux}, produit des plans de mitigation, adapte la communication
    \item Produit des \textbf{rapports d'avancement} et tableaux de bord pour les parties prenantes et la direction
    \item \textbf{Amélioration continue} : formations, ateliers de rétroaction, bilan des leçons apprises
    \item Automatise les tâches de gestion grâce à des agents IA pour gagner en efficacité opérationnelle
\end{itemize}

\subsection*{Intégrateur Sénior Qualité | Senior Quality Integrator} \hfill \textit{Octobre 2023 -- Novembre 2025}\\
\textbf{Banque Nationale du Canada (BNC)} -- Montréal, QC
\begin{itemize}
    \item Coordonne des équipes qualité multi-sites (Thaïlande) ; \textbf{gestion de plusieurs projets} simultanés
    \item Met en place des métriques et tableaux de bord (Datadog, Splunk, Jira) pour le suivi de performance
    \item Recrute, onboarde et forme en continu les membres de l'équipe
\end{itemize}

\subsection*{Intégrateur Qualité -- SDET | Quality Integrator -- Software Development Engineer in Test} \hfill \textit{Janvier 2022 -- Octobre 2023}\\
\textbf{Banque Nationale du Canada (BNC)} via Groupe Astek -- Montréal, QC
\begin{itemize}
    \item Conception et suivi de la stratégie de test ; \textbf{documentation projet} et métriques qualité
    \item Stack : Selenium, AppliTools, RestAssured, JMeter, Gatling, K6
\end{itemize}

\subsection*{QA Automaticien | QA Automation Engineer} \hfill \textit{Juin 2021 -- Décembre 2021}\\
\textbf{AODocs} -- Paris, France
\begin{itemize}
    \item Automatisation des tests non-régression ; rapports Allure ; POC Flutter
\end{itemize}

\subsection*{Automaticien CI/CD | DevOps Automation Engineer} \hfill \textit{Janvier 2020 -- Juin 2021}\\
\textbf{ENGIE DIGITAL} via Groupe Hénix -- Paris, France
\begin{itemize}
    \item \textbf{Pilotage projets} via Jira : administration, suivi d'avancement, mise en place de KPIs
    \item Création de la chaîne CI/CD : Jenkins, Docker, GitHub
\end{itemize}

\subsection*{Product Owner} \hfill \textit{Juin 2019 -- Septembre 2019}\\
\textbf{Hénix} -- Montrouge, France
\begin{itemize}
    \item Gestion du backlog, priorisation, animation de cérémonies, support et documentation clients
\end{itemize}

\subsection*{Développeur Logiciel | Software Developer} \hfill \textit{Mai 2017 -- Mai 2019}\\
\textbf{MEDIAPOST} via Groupe Hénix -- Paris, France
\begin{itemize}
    \item Développement back office, webservices, scripts BDD ; \textbf{reporting} et macros Excel
\end{itemize}

\subsection*{Consultant QA \& Automatisation | QA Automation Consultant} \hfill \textit{Avril 2016 -- Avril 2017}\\
\textbf{ENGIE IT} via Groupe Hénix -- Saint-Ouen, France
\begin{itemize}
    \item Pilotage opérationnel des environnements, gestion des incidents, administration plateforme
\end{itemize}

\section*{Certifications \& Formations}

\textbf{PMP Certification} (en cours) \hfill \textit{2026}

Project Management Professional -- Project Management Institute (PMI)

\textbf{AWS Cloud Practitioner} \hfill \textit{2026}

Amazon Web Services (AWS)

\textbf{Certified Scrum Master (CSM)} \hfill \textit{2024}

Scrum Alliance

\textbf{Jira ACP-600 -- Project Administration in Jira Server} \hfill \textit{2019}

Atlassian Certified Professional

\textbf{Cursus Automaticien et DevOps} \hfill \textit{2019}\\
\textit{AFCEPF -- Montrouge, France}

\textbf{Cursus Architecture logiciel \& Analyste informaticien - Certifié Master II} \hfill \textit{2017}\\
\textit{AFCEPF-HENIX -- Malakoff, France}

\section*{Formation Académique}

\textbf{Licence en Gestion (Bachelor's Degree in Management)} \hfill \textit{2011}\\
École Supérieure de Gestion (E.S.G) -- Paris, France

\section*{Compétences Linguistiques}

\textbf{Français:} Langue maternelle (Native)\\
\textbf{Anglais:} Niveau Avancé (Advanced / Professional Working Proficiency)

\end{document}
